\chapter*{Bibliographie et webographie}
\addcontentsline{toc}{chapter}{Bibliographie et webographie}

\begin{thebibliography}{99}

\bibitem{notepedagogique}
Institution de rattachement, \textit{Rapport de stage - note pedagogique}, document de reference remis pour l'encadrement du memoire de stage.

\bibitem{springboot}
Spring, \textit{Spring Boot Reference Documentation}, documentation officielle de Spring Boot.

\bibitem{angular}
Angular, \textit{Angular Documentation}, documentation officielle du framework Angular.

\bibitem{camunda}
Camunda, \textit{Camunda 7 Documentation}, documentation officielle du moteur BPMN Camunda.

\bibitem{keycloak}
Keycloak, \textit{Server Administration Guide}, documentation officielle de Keycloak.

\bibitem{alfresco}
Alfresco, \textit{Alfresco Content Services Documentation}, documentation officielle de la GED Alfresco.

\bibitem{sonarqube}
SonarSource, \textit{SonarQube Documentation}, documentation officielle relative a l'analyse statique et a la qualite logicielle.

\bibitem{argocd}
Argo Project, \textit{Argo CD Documentation}, documentation officielle relative a la synchronisation GitOps.

\bibitem{kubernetes}
Kubernetes, \textit{Kubernetes Documentation}, documentation officielle du socle d'orchestration conteneurisee.

\bibitem{ollama}
Ollama, \textit{Ollama Documentation}, documentation officielle du runtime d'execution locale des modeles IA.

\bibitem{supportflowdocs}
Projet SupportFlow, \textit{Documentation interne du depot : README, guides fonctionnels, guides DevOps/GitOps, guide SonarQube et support de soutenance}.

\end{thebibliography}
